\chapter{向量代数与空间解析几何}
\section{向量、距离与夹角公式}
\[
\bm a\cdot\bm b=|\bm a||\bm b|\cos\theta=\sum a_ib_i,\qquad
|\bm a\times\bm b|=|\bm a||\bm b|\sin\theta.
\]
叉积用三阶行列式，方向遵循右手法则；标量三重积 $[\bm a,\bm b,\bm c]=(\bm a\times\bm b)\cdot\bm c$ 是有向体积，取绝对值得平行六面体体积，除以 6 得四面体体积。非零向量垂直等价于点积为零；平行等价于叉积为零。投影长度为 $\bm a\cdot\bm b/|\bm b|$，投影向量为 $(\bm a\cdot\bm b)\bm b/|\bm b|^2$。

\ruleitem{直线与平面}{$\bm r=\bm r_0+t\bm v$；平面 $\bm n\cdot(\bm r-\bm r_0)=0$。方向向量某分量为零时，不应在对称式中除以该分量，应写对应坐标固定。两个平面相交线方向为 $\bm n_1\times\bm n_2$。}
\ruleitem{夹角}{两直线取锐角时 $\cos\theta=|\bm v_1\cdot\bm v_2|/(|\bm v_1||\bm v_2|)$；两平面用法向量同式；线面夹角 $\sin\theta=|\bm v\cdot\bm n|/(|\bm v||\bm n|)$。}
\ruleitem{距离}{
\[
d(P,\Pi)=\frac{|Ax_0+By_0+Cz_0+D|}{\sqrt{A^2+B^2+C^2}},\quad
d(P,l)=\frac{|(\bm P-\bm P_0)\times\bm v|}{|\bm v|}.
\]
两不平行直线间最短距离为
$\displaystyle d=\frac{|(\bm P_2-\bm P_1)\cdot(\bm v_1\times\bm v_2)|}{|\bm v_1\times\bm v_2|}$。
平行时改用点到直线距离，不可把分母为零的式子硬算。}
\ruleitem{平面束}{过两相交平面 $L_1=0,L_2=0$ 的交线的平面为 $\lambda L_1+\mu L_2=0$，$(\lambda,\mu)\ne(0,0)$。写成 $L_1+kL_2=0$ 会遗漏 $L_2=0$ 本身。}

\section{常见曲面与切平面}
\begin{longtable}{p{.45\textwidth}p{.44\textwidth}}
\toprule 标准方程（$a,b,c>0$） & 几何含义\\\midrule\endhead
$x^2/a^2+y^2/b^2+z^2/c^2=1$ & 椭球面\\
$x^2/a^2+y^2/b^2-z^2/c^2=1$ & 单叶双曲面，轴为 $z$ 轴\\
$z^2/c^2-x^2/a^2-y^2/b^2=1$ & 双叶双曲面，轴为 $z$ 轴\\
$z=x^2/a^2+y^2/b^2$ & 椭圆抛物面\\
$z=x^2/a^2-y^2/b^2$ & 双曲抛物面（马鞍面）\\
$x^2/a^2+y^2/b^2-z^2/c^2=0$ & 椭圆锥面\\
缺少某坐标的曲面方程 & 通常为沿该坐标方向延伸的柱面\\\bottomrule
\end{longtable}
通过平移完成平方，把非标准方程化为标准形式；截面法有助于识别内部区域。旋转曲面绕 $z$ 轴时，常把平面曲线的径向变量换为 $\sqrt{x^2+y^2}$。

对光滑曲面 $F(x,y,z)=0$，若 $\grad F(P)\ne0$，切平面为 $\grad F(P)\cdot(\bm r-\bm P)=0$，法线方向为 $\grad F(P)$。图像型 $z=f(x,y)$ 可取法向 $(f_x,f_y,-1)$。两曲面交线的切向为两法向叉积，要求叉积非零。

\section{空间几何真题}
\problem{48}{异面直线距离}{2010 年第 2 届，非数学类，一(5)，\src{02}}
$l_1:x=y,z=0$，$l_2:(x-2)/4=(y-1)/(-2)=(z-3)/(-1)$，求距离。
\sol 取 $P_1=(0,0,0)$、$P_2=(2,1,3)$，$\bm v_1=(1,1,0)$、$\bm v_2=(4,-2,-1)$，有 $\bm v_1\times\bm v_2=(-1,1,-6)$。所以
\[d=\frac{|(2,1,3)\cdot(-1,1,-6)|}{\sqrt{38}}=\ans{\frac{19}{\sqrt{38}}=\frac{\sqrt{38}}2}.\]
取绝对值后距离才非负，不能直接用有向三重积作答案。

\problem{49}{过交线的两个垂直平面}{2012 年第 4 届，非数学类，一(2)，\src{04}}
两平面都过 $2x+y-3z+2=0$ 与 $5x+5y-4z+3=0$ 的交线，且互相垂直，其中一个过 $(4,-3,1)$。求二平面。
\sol 该点代入两原平面左端均为 4，故过该点的平面束参数满足 $\lambda+\mu=0$，得到 $\Pi_1:3x+4y-z+1=0$。
另一个的法向为 $(2\lambda+5\mu,\lambda+5\mu,-3\lambda-4\mu)$，与 $(3,4,-1)$ 点积为 $13\lambda+39\mu=0$，取 $\lambda=-3\mu$，得
\[\ans{\Pi_1:3x+4y-z+1=0,\quad\Pi_2:x-2y-5z+3=0}.\]

\problem{50}{指定平行方向的切平面}{2016 年第 8 届，非数学类，一(5)，\src{08}}
求曲面 $z=x^2/2+y^2$ 上与平面 $2x+2y-z=0$ 平行的切平面。
\sol 法向 $(x,2y,-1)$ 与 $(2,2,-1)$ 平行，因第三分量相同，比例为 1，所以切点 $(2,1,3)$。切平面为 $2(x-2)+2(y-1)-(z-3)=0$，即 $\ans{2x+2y-z=3}$。

\problem{51}{过两点且垂直另一平面}{2025 年第 17 届，非数学 A/B 类，一(2)，共题，\src{17A}\src{17B}}
平面过原点及 $(6,-3,2)$，并垂直于 $4x-y+2z-8=0$，求方程。
\sol 所求法向同时垂直于 $(6,-3,2)$ 与 $(4,-1,2)$，可取它们的叉积 $(-4,-4,6)$，约成 $(2,2,-3)$，所以 $\ans{2x+2y-3z=0}$。常数 $-8$ 影响参照平面位置，不影响它的法向。
